%% An explicit equational basis presenting every one-sorted Lawvere theory.

\subsection{An equational presentation of every Lawvere theory}
\label{sec:p10-lawvere-presentations}

Let \(\mathcal T\) be a Lawvere theory in the sense of
Definition~\ref{def:p10-lawvere-theory}.
Write \(p_i^n:n\to1\), \(0\leq i<n\), for its selected projections.
The structural functor fixes \(p_0^1=1_1\).
For arrows \(g_i:n\to1\), write
\(\langle g_0,\ldots,g_{k-1}\rangle:n\to k\) for their unique
tuple.  An empty tuple is the unique arrow \(!_n:n\to0\).
Let \(X_n=(x_0,\ldots,x_{n-1})\), and abbreviate this variable
list by \(\mathbf x_n\).

\begin{hbdefinition}{def:p10-lawvere-named-presentation}{The named operation presentation}
The signature \(\Sigma_{\mathcal T}\) has one \(n\)-ary symbol
\(\sigma_f\) for every arrow \(f:n\to1\) in \(\mathcal T\),
for every \(n\geq0\).
Its equation set \(E_{\mathcal T}\), in the explicit finite contexts
of Definition~\ref{def:p10-lawvere-presentation}, consists of these
two families:
\[
\begin{aligned}
 \sigma_{p_i^n}(\mathbf x_n)&\approx x_i
       &&\text{in }X_n,\quad 0\leq i<n,\\
 \sigma_f\bigl(\sigma_{g_0}(\mathbf x_n),\ldots,
                    \sigma_{g_{k-1}}(\mathbf x_n)\bigr)
   &\approx \sigma_{\,f\circ\langle g_0,\ldots,g_{k-1}\rangle}
                      (\mathbf x_n)
       &&\text{in }X_n.
\end{aligned}
\]
The second family ranges over \(n,k\geq0\), \(f:k\to1\), and
all \(k\)-tuples \(g_i:n\to1\).
At \(k=0\) it reads
\[
                  \sigma_f()\approx\sigma_{f\circ!_n}(\mathbf x_n)
                        \quad\text{in }X_n,\qquad f:0\to1.
\]
Thus it also specifies how a nullary term is regarded in a
context with unused variables.  There are no projection equations
in \(X_0\).
Smallness of \(\mathcal T\) makes the signature and these indexed
equation families small sets.

For a term \(t\) in \(X_n\), define its interpretation
\(\llbracket t\rrbracket_n:n\to1\) recursively by
\[
 \llbracket x_i\rrbracket_n=p_i^n,\qquad
 \llbracket\sigma_f(t_0,\ldots,t_{k-1})\rrbracket_n
       =f\circ\langle\llbracket t_0\rrbracket_n,\ldots,
                              \llbracket t_{k-1}\rrbracket_n\rangle .
\]
The empty tuple gives
\(\llbracket\sigma_f()\rrbracket_n=f\circ!_n\).
In this section abbreviate the relation
\(\equiv_{E_{\mathcal T}}^n\) of
Definition~\ref{def:p10-lawvere-presentation} by
\(\equiv_{E_{\mathcal T},n}\).
\index{Lawvere theory!named presentation}
\end{hbdefinition}

\begin{hbtheorem}{thm:p10-lawvere-presentation-recovery}{proved=annex:p10.lawvere-present.1}{Recovery of an abstract theory from its named equations}
Every term \(t\) in \(X_n\) satisfies
\[
       t\equiv_{E_{\mathcal T},n}
                         \sigma_{\llbracket t\rrbracket_n}(\mathbf x_n).
\]
Consequently, for any terms \(s,t\) in the same context,
\[
 s\equiv_{E_{\mathcal T},n}t
       \quad\Longleftrightarrow\quad
                  \llbracket s\rrbracket_n=\llbracket t\rrbracket_n.
\]
Interpretation induces an isomorphism of categories
\[
 \varepsilon_{\mathcal T}:
       T_{\Sigma_{\mathcal T},E_{\mathcal T}}
                  \longrightarrow\mathcal T
\]
which is the identity on objects and preserves the selected
finite products and the structural finite-set maps.
For a tuple arrow its value is
\[
       \varepsilon_{\mathcal T}([t_0],\ldots,[t_{m-1}])
          =\langle\llbracket t_0\rrbracket_n,\ldots,
                                     \llbracket t_{m-1}\rrbracket_n\rangle.
\]
The inverse sends \(h:n\to m\) to
\[
       \bigl([\sigma_{p_j^m\circ h}(\mathbf x_n)]\bigr)_{0\leq j<m}.
\]
For \(m=0\) this is the empty tuple.
All assertions include \(n=0\), nullary symbols, and composition
through the terminal object \(0\).  Faithfulness of the structural
finite-set functor is not required.
\end{hbtheorem}

The named signature and quotient of its term theory by evaluation
are given in
\cite[Definition~13.5 and Proposition~13.7, p.~123;
Remark~13.8, p.~124]{AdamekRosickyVitale2010Draft}.
The proof here identifies that evaluation kernel using the two
explicit equation families above.
Lawvere's thesis treats presentations and substitution in
\cite[Chapter~II, \S2, pp.~69--73]{Lawvere2004FunctorialReprint}.
Its theory category uses finite coproducts and operations \(1\to n\);
the product convention used here takes its opposite
\cite[Chapter~II, \S1, pp.~61--63]{Lawvere2004FunctorialReprint}.
The latter locators are pages of the authorized 2004 reprint;
the former are pages of the February 2010 author draft.

\begin{hbdefinition}{def:p10-lawvere-theory-morphism}{Morphisms of one-sorted theories}
A theory morphism \(H:\mathcal T\to\mathcal U\) is a functor
which is the identity on the objects \(n\) and satisfies
\[
                         H(p_i^{n,\mathcal T})=p_i^{n,\mathcal U}.
\]
Equivalently, it commutes with the structural functors from the
opposite of the finite-ordinal category.
Composition and identity functors give the category of these
theories and morphisms.
Write \(\operatorname{Mod}(\mathcal T)\) for the models and all
natural transformations of Definition~\ref{def:p10-lawvere-model}.
This morphism convention is
\cite[Definition~11.3 and Remark~11.4, p.~98]{AdamekRosickyVitale2010Draft}.
\end{hbdefinition}

\begin{hbtheorem}{thm:p10-lawvere-model-restriction}{proved=annex:p10.lawvere-present.2}{Theory morphisms and restriction of models}
Every theory morphism \(H:\mathcal T\to\mathcal U\) preserves
the selected finite products.  It induces a functor
\[
 H^*:\operatorname{Mod}(\mathcal U)\longrightarrow
                         \operatorname{Mod}(\mathcal T),
 \qquad H^*(M)=M\circ H,\qquad
              (H^*\alpha)_n=\alpha_n.
\]
This restricts to the full subcategories with nonempty carrier
\(M(1)\).  For \(K:\mathcal U\to\mathcal V\),
\[
             (K\circ H)^*=H^*\circ K^*,\qquad (1_{\mathcal T})^*=1.
\]
On named presentations, replacing \(\sigma_f\) by \(\sigma_{H(f)}\)
and fixing variables induces a theory morphism \(\widehat H\)
between the corresponding term theories, satisfying
\[
       \varepsilon_{\mathcal U}\circ\widehat H
                        =H\circ\varepsilon_{\mathcal T}.
\]
\end{hbtheorem}
